%% LyX 2.4.4 created this file.  For more info, see https://www.lyx.org/.
%% Do not edit unless you really know what you are doing.
\documentclass[oneside,english]{amsart}
\usepackage[T1]{fontenc}
\usepackage[latin9]{inputenc}
\synctex=-1
\usepackage{mathrsfs}
\usepackage{amsthm}

\makeatletter
%%%%%%%%%%%%%%%%%%%%%%%%%%%%%% Textclass specific LaTeX commands.
\newlength{\lyxlabelwidth}      % auxiliary length 

%%%%%%%%%%%%%%%%%%%%%%%%%%%%%% User specified LaTeX commands.
\date{}
\usepackage{commath}

\makeatother

\usepackage{babel}
\begin{document}
{[}{[}At times I found myself saying ``dotted red line'' a lot.
Your thoughts on introducing the word ``segment'' for any line segment
appearing on a tile?{]}{]}

{[}{[}The following is a replacement for Section~2. At the beginning
I note changes from what's there. Then I just go off on my own.{]}{]}

For positive integers $m$ and {[}{[}word ``and'' added here{]}{]}
$n$, an \emph{$\left(m,n\right)$ binary lattice} is a rectangular
lattice of $m+1$ by $n+1$ vertices, with each vertex labeled $0$
or $1$. We also define a \emph{framed lattice} to be a binary lattice
in which all {[}{[}``the'' changed to ``all'' here{]}{]} boundary
vertices are labeled $0$. An example of a $\left(5,7\right)$ framed
lattice is shown on the right of Figure~{[}{[}ref 3{]}{]}. {[}{[}Next
stuff reworded. E.g., Definition 2.3 now in text, not a definition
environment.{]}{]} A \emph{cell} is a $\left(1,1\right)$ binary lattice,
{[}{[}the next text is reworded{]}{]} which we generally draw without
connecting lines. The sixteen possible cells are pictured in Figure~{[}{[}make
existing text into a Figure to be referenced{]}{]}. A binary lattice
or a framed lattice is \emph{proper} if it contains no cells $\begin{array}{cc}
0 & 1\\
1 & 0
\end{array}$ or $\begin{array}{cc}
1 & 0\\
0 & 1
\end{array}$. {[}{[}Note: unless we reword this definition, we should never say
``proper lattice'' but rather ``proper binary lattice.'' I prefer
``proper'' to be a modifier, so that ``framed'' replaces ``binary''
but ``proper'' doesn't. I know this is a change from my previous
preference. I never claimed consistency.{]}{]}

{[}{[}Based on this wording I think $\left(m,n\right)$ goes between
``proper'' and ``framed.'' If you want something else, e.g., $\left(m,n\right)$
proper framed lattices, then let's talk about how to write the definitions
accordingly.{]}{]}

{[}{[}In April~8 draft the following was Definition~2.3. I'm thinking
it might make more sense to just have this in text and not a Definition
environment. Thoughts?{]}{]} 

{[}{[}Current Definition~2.3 uses ``vertex'' for a mosaic which
I think might be confusing. Thus I reworded.{]}{]}

{[}{[}Throughout the paper, I recommend ${\tt \backslash mathscr}$
for $\mathscr{L}$ and ${\tt \backslash mathcal}$ for $\mathcal{M}$
so that they look quite different.{]}{]} Let $\mathcal{M}$ be a polygon
mosaic. We will create a binary lattice from $\mathcal{M}$ by replacing
every point that is at a corner of a tile of $\mathcal{M}$ by either
$0$ or $1$. For each such point, we count the number of polygons
of $\mathcal{M}$ that surround it, i.e., polygons for which the point
is in the polygon's interior. If this number is even, label the vertex
$0$, and if it is odd, label the vertex $1$. We then delete all
red dotted lines that are part of polygons of $\mathcal{M}$. We define
$\mathscr{L}\left(\mathcal{M}\right)$ to be the resulting binary
lattice. See Figure~{[}ref for Fig~3{]}. {[}{[}maybe the caption
of Fig~3 should change from ``$\mathscr{L}$ applied to a polygon
mosaic'' to ``Applying $\mathscr{L}$ to a polygon mosaic''??{]}{]}
If $\mathcal{M}$ is an $\left(m,n\right)$ polygon matrix, i.e.,
$m$ tiles high and $n$ tiles wide, then $\mathscr{L}\left(\mathcal{M}\right)$
is an $\left(m,n\right)$ binary lattice, i.e., $m+1$ vertices high
and $n+1$ vertices wide. In fact, $\mathscr{L}\left(\mathcal{M}\right)$
is an $\left(m,n\right)$ framed lattice because every point on the
edge of $\mathcal{M}$ is surrounded by no polygons.

In our definition of $\mathscr{L}$, a tile of a polygon mosaic $\mathcal{M}$
corresponds to a cell of $\mathscr{L}\left(\mathcal{M}\right)$. For
tile of $\mathcal{M}$ that is not the blank tile, consider laying
its corresponding cell over it. The red dotted line of the tile separates
the vertices of the cells, with vertices labeled $0$ on one side
and labeled $1$ on the other. For example, the second tile of Figure~{[}ref
1{]} must map to either the cell $\begin{array}{cc}
0 & 0\\
1 & 0
\end{array}$ or the cell $\begin{array}{cc}
1 & 1\\
0 & 1
\end{array}$. Which one? Consider first where the interior of this polygon is
relative to this dotted red line, and second, the number of other
polygons that contain this polygon. For example, the tile maps to
$\begin{array}{cc}
0 & 0\\
1 & 0
\end{array}$ if the interior is below the dotted red line and an even number of
polygons contain this polygon. On the other hand, a blank tile of
$\mathcal{M}$ corresponds to a cell that is either $\begin{array}{cc}
0 & 0\\
0 & 0
\end{array}$ or $\begin{array}{cc}
1 & 1\\
1 & 1
\end{array}$. With seven tiles, and no cell able to be the image of distinct tiles,
$\mathscr{L}\left(\mathcal{M}\right)$ can contain fourteen types
of cells. The two types that never appear in $\mathscr{L}\left(\mathcal{M}\right)$
are $\begin{array}{cc}
0 & 1\\
1 & 0
\end{array}$ and $\begin{array}{cc}
1 & 0\\
0 & 1
\end{array}$. That is, $\mathscr{L}\left(\mathcal{M}\right)$ is proper.

Given any proper binary lattice $\ell$, we can reverse this process,
creating a mosaic by replacing each cell by the unique tile that can
map to it. Define $\mathscr{L}^{-1}$ by restricting this map to proper
framed lattices, and let us show that $\mathscr{L}^{-1}$ is indeed
an inverse to $\mathscr{L}$, thus proving that $\mathscr{L}$ gives
a bijection from $\left(m,n\right)$ polygon mosaics to proper $\left(m,n\right)$
framed lattices. Our discussion above shows that $\mathscr{L}^{-1}\circ\mathscr{L}\left(\mathcal{M}\right)=\mathcal{M}$
for any polygon mosaic $\mathcal{M}$, so fixing a proper framed lattice
$\ell$, we must show that $\mathscr{L}^{-1}\left(\ell\right)$ is
a polygon mosaic and that $\mathscr{L}\circ\mathscr{L}^{-1}\left(\ell\right)=\ell$.

First observe that polygon mosaics can be characterized as those mosaics
in which every end of a dotted red line on a tile meets, on an adjacent
tile, another end of a dotted red line. In the mosaic $\mathscr{L}^{-1}\left(\ell\right)$,
these ends are exactly the points which, in the corresponding cell
of $\ell$, lie halfway between adjacent vertices of $\ell$, one
labeled $0$ and the other labeled $1$. Because $\ell$ is a framed
lattice, the vertex labeled $1$ is an internal vertex. Thus, another
cell of $\ell$ shares these two vertices, and the dotted red lines
of $\mathscr{L}^{-1}\left(\ell\right)$ from the two corresponding
tiles meet. Therefore $\mathscr{L}^{-1}\left(\ell\right)$ is a polygon
mosaic.

By our analysis of $\mathscr{L}$ and our definition of $\mathscr{L}^{-1}$,
each cell of $\mathscr{L}\circ\mathscr{L}^{-1}\left(\ell\right)$
must either be equal to the corresponding cell of $\ell$ or is formed
from the corresponding cell by relabeling all vertices $0$ as $1$
and $1$ as $0$. Furthermore, since the cells are connected, i.e.,
none is isolated, we must then have that $\mathscr{L}\circ\mathscr{L}^{-1}\left(\ell\right)$
itself is either equal to $\ell$ or formed by relabeling all vertices
$0$ as $1$ and $1$ as $0$. By hypothesis, $\ell$ is a framed
lattice, and by construction, every binary lattice in the range of
$\mathscr{L}$ is a framed lattice. Therefore $\mathscr{L}\circ\mathscr{L}^{-1}\left(\ell\right)=\ell$.
This concludes our proof that $\mathscr{L}^{-1}$ is the inverse of
$\mathscr{L}$ and thus that $\mathscr{L}$ is bijective.

\bigskip{}
\bigskip{}

{[}{[}To insert when defining $v$ with signs (at least something
along these lines){]}{]}

In Section~{[}ref{]}, we remarked that $\begin{array}{cc}
0 & 0\\
1 & 0
\end{array}$ and $\begin{array}{cc}
1 & 1\\
0 & 1
\end{array}$ map to the same tile under $\mathscr{L}^{-1}$. By contrast, it is
$\begin{array}{cc}
0 & 0\\
1 & 0
\end{array}$ and $\begin{array}{cc}
1 & 0\\
1 & 1
\end{array}$ whose image under $v$ is negative.\bigskip{}
\bigskip{}

{[}{[}Possible insertion after stating or after proving that ``For
a polygon mosaic $P$ and $\ell=\mathscr{L}\left(P\right)$, if $P$
contains an even {[}resp. odd{]} number of polygons then $v\left(\ell\right)$
is positive {[}resp. negative{]}''.{]}{]}

We might want to think of this result as relating a ``global'' characteristic
and ``local'' characteristic, since the number of polygons in $P$
is determined by looking at the entire mosaic, whereas $v\left(\ell\right)$
is calculated by looking at one cell at a time. The reader may wish
to compare this to the definition of $\mathscr{L}$, which is done
``globally'' and the definition of $\mathscr{L}^{-1}$, which was
defined ``locally.''
\end{document}
